\chapter{货币市场与短期流动性契约}
\label{chap:Money_Markets_and_Short_Term_Liquidity_Contracts}

\begin{openingcase}{当隔夜融资突然变得昂贵}

2013~年~6~月，中国银行间货币市场经历了异常剧烈的流动性紧张。此前通常处于较低个位数水平的短期回购利率突然上升，6~月~20~日，部分基准隔夜和七天期回购利率在日内达到异常高位。银行同业放贷更为谨慎，交易活动也随之收缩。

这次紧张首先出现在期限极短的银行间融资中，并非源于银行资产长期价值的明显恶化。季节性现金需求、税款缴纳、资本流动变化，以及央行流动性条件收紧，与短期利率上升同时出现。随着不确定性增加，银行同业放贷更加谨慎，交易活动减弱。

一个日常用于调剂流动性的市场，为什么会在如此短的时间内陷入严重紧张？

\end{openingcase}

\begin{tcolorbox}[
  enhanced,
  colback=FinanceBlue!6,
  colframe=FinanceBlue,
  boxrule=0.6pt,
  sharp corners,
  left=8pt,
  right=8pt,
  top=7pt,
  bottom=7pt,
  before skip=1em,
  after skip=1.6em,
  title=\textbf{核心问题},
  coltitle=white,
  colbacktitle=FinanceBlue,
  fonttitle=\bfseries,
  attach boxed title to top left={xshift=6pt,yshift=-2mm},
  boxed title style={
    boxrule=0pt,
    sharp corners,
    left=5pt,
    right=5pt,
    top=2pt,
    bottom=2pt
  }
]
\begin{enumerate}[
  leftmargin=*,
  label=\textcolor{FinanceBlue}{\textbf{\arabic*.}},
  itemsep=0.55em,
  topsep=0.2em
]

\item 即使基础经济活动持续较长时间，企业和金融机构为什么仍需要短期融资？

\item 无担保和有抵押货币市场契约如何重新分配流动性，并决定短期融资条款？

\item 短期融资条件如何定价？通常流动性充裕的货币市场为什么会在压力下变得脆弱？

\end{enumerate}
\end{tcolorbox}

\section{短期融资的经济逻辑}

短期融资使企业和金融机构无需持有大量闲置现金，也能满足支付需求。其作用在于提供流动性和灵活性，但较短期限也带来了反复再融资的需要。

\subsection{现金流时点与短期融资需求}

现金收入与现金支出并不总在同一时点发生。具有偿付能力的企业仍可能面临暂时性现金短缺。短期融资弥合这些时间缺口，减少参与者需要持有的闲置现金。

假设一家企业期初现金余额为人民币~20~万元，须在一周内向员工和供应商支付~80~万元，但预计在~45~天后从客户处收到~100~万元。\footnote{译注：原稿此处为~45~天，下文与图示均为第~28~天。译稿保留原数值，待作者核定。}该企业的现金状况变化如下：

\begin{itemize}
\item 期初，企业持有人民币~20~万元现金。

\item 第~7~天付款到期时，企业面临~60~万元现金短缺。

\item 第~28~天，客户支付~100~万元，企业得以偿还短期借款，并保留正的现金余额。
\end{itemize}

因此，即使该期间预期现金流入超过计划现金流出，企业仍需要~60~万元临时融资。这说明，收付款时点差异可能产生无法从期末现金余额看出的融资需求。

金融机构也存在类似的时点问题。银行可能在款项流入前先发生结算资金流出；证券交易商可能需要为持有较长时间的证券库存续借隔夜融资。保证金支付和集中的支付义务也可能产生暂时性融资需求。

短期现金短缺并不必然意味着资不抵债。偿付能力关注资产经济价值是否足以覆盖负债，流动性则关注能否在支付到期时履行义务。因此，一家机构即使具有偿付能力，若无法及时将资产变现或取得新融资以履行支付义务，也可能面临流动性短缺。
\begin{figure}[H]
\centering
\begin{tikzpicture}[
  >=Latex,
  axis/.style={-Latex, line width=0.8pt},
  inflow/.style={-Latex, line width=1.0pt, FinanceBlue},
  outflow/.style={-Latex, line width=1.0pt}
]

\draw[axis] (0,0) -- (10.5,0) node[right]{时间};

% Tick labels
\node[below=4pt] at (0,0) {0};
\node[below=4pt] at (2.2,0) {7};
\node[below=4pt] at (8,0) {28};

% Initial cash
\draw[inflow] (0,0) -- (0,0.5)
  node[above=2pt, align=center]
   {\small 期初现金余额\\\small 人民币~20~万元};
% Cash outflow
\draw[outflow] (2,0) -- (2,-1.5)
  node[below=2pt, align=center]
  {\small 向供应商支付\\\small 及向员工支付\\\small 人民币~80~万元};

% Cash inflow
\draw[inflow] (8,0) -- (8,1.6)
  node[above=2pt, align=center]
  {\small 客户应收账款回款\\\small 人民币~100~万元};

\end{tikzpicture}

\caption{现金流时点缺口与短期融资的作用}
\label{fig:ch7-cash-flow-timing}
\end{figure}

金融体系内部也存在类似的时点问题。一家机构当天收到的现金可能超过其需要，另一家则面临暂时性缺口。货币市场重新分配这些资金余额，使参与者能够节约闲置流动性，并支持支付和证券交易结算。

融资安排的期限影响风险。如果时点缺口确属暂时性，短期借款可以较好地匹配实际需要。但如果短期债务用于融资的资产多年内都缺乏流动性，借款人就必须反复再融资。这种期限错配产生滚动融资风险，即现有融资到期时，替代融资变得更昂贵或无法获得的风险。

\begin{example}{管理多项现金流义务}
{multiple-cash-flow-obligations}

一家制造企业在一个~15~天期间开始时持有人民币~200~万元现金。期间，企业面临多项支付义务，并预计在不同日期收到客户付款。新增融资前的预计现金流如下。

\begin{center}
\begin{tabular}{lrr}
\toprule
\textbf{日期} &
\shortstack[c]{\textbf{现金流}\\\textbf{（人民币百万元）}} &
\shortstack[c]{\textbf{累计}\\\textbf{现金余额}} \\
\midrule
期初 & 2.0 & 2.0 \\
第~2~天：供应商付款 & -5.0 & -3.0 \\
第~4~天：客户回款 & 1.5 & -1.5 \\
第~5~天：工资和税款 & -1.0 & -2.5 \\
第~8~天：客户回款 & 4.0 & 1.5 \\
第~10~天：运费和利息 & -0.5 & 1.0 \\
第~15~天：客户回款 & 2.0 & 3.0 \\
\bottomrule
\end{tabular}
\end{center}

企业在第~2~天出现最大现金缺口，预计现金余额降至负~300~万元。因此，企业至少需要~300~万元临时融资，才能履行支付义务。此后的客户付款使企业能够减少或偿还借款。

这一例子说明，短期融资需求取决于现金流的先后顺序，而不仅是给定期间内总收入与总支出的差额。虽然企业期末现金余额为正，但若无临时性外部融资，就无法按时履行全部义务。

承诺信贷额度或短期贷款可以弥合这一缺口。融资选择取决于可得性、成本、期限及预期客户付款的可靠性。

\end{example}

\subsection{营运资本、流动性管理与财务灵活性}

对非金融企业而言，短期融资与营运资本密切相关。在现金回收之前，经营活动会使资金占用于存货和应收账款，而供应商提供的商业信用则是短期融资来源。

经营性营运资本的一种简化计量为

\begin{equation*}
OWC=\text{存货}+\text{应收账款}-\text{应付账款}.
\label{eq:ch7-operating-working-capital}
\end{equation*}

存货或应收账款增加会占用资金，应付账款增加则提供短期融资。如果必须在客户付款之前投入资金，即使企业盈利，销售快速增长也可能增加短期融资需求。

流动性管理涉及用于满足短期现金需求的资产和融资安排。企业可以持有现金和有价证券、维持银行承诺授信，或取得短期市场融资。这些选择反映流动性保障与融资成本之间的权衡。更大的流动性缓冲可降低暂时性融资中断的影响，但流动资产的回报通常低于生产性投资。承诺信贷额度以一定成本保障资金获取，而依赖极短期市场融资则会增加滚动融资风险暴露。

财务灵活性指在不承担高额成本、不干扰投资与经营的情况下，获得融资或调整融资安排的能力。它不仅取决于当前现金持有量，还取决于未使用的借款能力、进入其他融资市场的渠道、未设权利负担的抵押品、负债期限结构和借款人信用状况。这些因素决定条件变化时能够筹集或重组融资的程度。

对金融机构而言，流动性管理和财务灵活性高度依赖资产构成及融资市场准入。银行可以持有可出售或质押的高质量流动性资产，在银行间市场融资，并以合格抵押品取得央行流动性。证券交易商通常通过回购协议融资证券库存，其融资能力因而部分取决于抵押品资格和折扣率。货币市场基金持有短期资产，但若这些资产无法在不作重大价格让步的情况下迅速出售，快速赎回仍可能带来流动性压力。

因此，短期融资在提供灵活性的同时，也产生对再融资的依赖。较短期限使融资条件能够迅速调整，但也要求借款人更频繁地回到贷款人或市场寻求资金。期限不仅决定偿还时点，也决定融资重新定价、续借或撤回的频率。

\begin{tcolorbox}[
    enhanced,
    breakable,
    colback=white,
    colframe=FinanceBlue,
    boxrule=0.7pt,
    sharp corners,
    borderline west={2.2pt}{0pt}{FinanceBlue},
    title=\textbf{金融视角：短期限与滚动融资依赖},
    coltitle=white,
    colbacktitle=FinanceBlue,
    fonttitle=\bfseries,
    before skip=1.2em,
    after skip=1.2em
]

短期借款允许融资条款频繁调整，从而提供灵活性，但同一特征也增加了再融资依赖。一笔一天期贷款的贷款人只承诺将资金提供至到期日，可在下次续借时重新定价或拒绝续借。如果基础融资需求仍然存在，借款人就继续面临市场条件变化。短期限因而提高了续借决策的重要性，并使融资风险集中于再融资时点。

\end{tcolorbox}

\section{货币市场契约与工具}

货币市场工具和契约是用于满足临时融资需求、管理现金余额及融资证券头寸的短期请求权与融资安排。它们在发行人或交易对手、抵押品、信用风险暴露、可交易性和报价惯例方面存在差异，进而影响定价、流动性和风险。

\subsection{国库券与短期主权请求权}

国库券是中央政府发行的短期债务证券，通常不定期支付票息，而在到期时支付固定面值，回报来自购买价格与到期价值之差。

若以~\(P\)~购买面值为~\(F\)~的国库券，持有~\(d\)~天至到期，其持有期回报率为

\begin{equation}
HPR=\frac{F-P}{P}.
\label{eq:ch7-tbill-hpr}
\end{equation}

为比较不同期限国库券，可将持有期回报率按单利年化。所得收益率取决于市场报价惯例。以购买价格为基础的货币市场收益率可写为

\begin{equation}
y_m=\frac{F-P}{P}\frac{D}{d},
\label{eq:ch7-money-market-yield}
\end{equation}
其中，$D$~为市场惯例规定的年度计息天数，通常为~360~或~365。

一些国库券市场还使用以面值为基础的贴现收益率：

\begin{equation}
y_d=\frac{F-P}{F}\frac{D}{d}.
\label{eq:ch7-discount-yield}
\end{equation}

\begin{example}{报价惯例改变报告收益率}{tbill-quotation-convention}

一张期限为~91~天、面值为~100~的国库券，售价为~98.80。采用每年~360~天的计息基础，持有期回报率为

\[
HPR=\frac{100-98.80}{98.80}=1.2146\%.
\]

年化货币市场收益率为

\[
y_m=\frac{1.20}{98.80}\frac{360}{91}=4.805\%,
\]
贴现收益率则为

\[
y_d=\frac{1.20}{100}\frac{360}{91}=4.747\%.
\]

两种收益率对应相同的基础现金流，数值差异完全来自报价惯例。

\end{example}

国库券是短期主权债务最常见的形式，但具体工具、期限和发行惯例因国家而异。

短期主权证券的作用不限于政府融资。当信用风险较低、二级市场流动性较高时，它们可以作为临时现金投资工具、有抵押融资的抵押品，以及其他短期请求权定价的基准。

信用风险较低，并不意味着所有持有期内市场价格都无风险。如果到期前出售国库券，其价格可能随市场利率变化。结算、市场流动性及主权信用条件也可能产生影响。与其他特征相近的较长期固定利率证券相比，短期限通常使其利率敏感性较低。

\subsection{商业票据与企业短期融资}

商业票据（commercial paper，CP）是企业和金融机构为满足短期融资需求而发行的、通常无担保的短期债务工具。投资者主要依赖发行人的信用状况和流动性状况。商业票据可采用贴现或付息方式发行。贴现票据低于面值发行，到期按面值偿还；付息票据则在本金之外支付约定利息。

两种结构的区别在于投资者获得回报的方式。贴现票据的回报来自购买价格与面值之差，付息票据的回报来自约定利息支付。

\begin{example}{贴现商业票据与付息商业票据}
{commercial-paper-payment-structures}

考虑两笔~90~天期商业票据，其本金或面值均为人民币~1000~万元。假设计息年度为~360~天。

\textbf{贴现商业票据。}

假设票据以~985~万元发行，90~天后按面值~1000~万元偿还，投资者获得~15~万元。年化贴现率为

\[
d=\frac{10-9.85}{10}\times\frac{360}{90}
=6\%.
\]

以购买价格为分母计算的相应单利年化投资收益率为

\[
y=\frac{10-9.85}{9.85}\times\frac{360}{90}
\approx 6.09\%.
\]

\textbf{付息商业票据。}

假设票据按面值发行，年利率为~6\%，到期应付利息为

\[
I=10\times 6\%\times\frac{90}{360}
=0.15\text{ 百万元}.
\]

投资者到期收到本息合计人民币~1015~万元，单利年化收益率为~6\%。

这一例子说明，贴现率与投资收益率不同。贴现率以面值为基础计算，投资收益率则以购买价格为基础计算。按面值发行的付息票据通过约定利率表达其回报。

虽然两种报价均为~6\%，但贴现票据的贴现率以面值而非购买价格计算，因而提供略高的投资收益率。

\end{example}

投资者要求对货币时间价值，以及与发行人相关的信用和流动性风险给予补偿。因此，必要收益率反映可比期限的基准利率，以及与发行人特定信用和流动性条件相关的利差。

对信用质量较高的发行人而言，商业票据提供直接获取短期市场融资的渠道，其成本可能低于银行借款。不过，依赖反复发行会产生滚动融资风险。依赖~30~天期票据的企业，必须在每次到期时以可用现金偿还，或获得新融资。如果市场对其信用质量的评价恶化，或投资者需求减弱，利差可能扩大，新发行渠道也可能受限。因此，企业常以承诺银行信贷额度或流动性储备支持商业票据计划。即使发行人仍有偿付能力，滚动融资风险也可能出现。

相关的短期市场融资形式是资产支持商业票据（ABCP）。与普通企业商业票据不同，ABCP~由特殊目的载体发行，以应收账款或贷款等金融资产池为支持。该计划还可能获得发起人或其他机构提供的流动性便利或信用增级。

\subsection{存单与银行融资}

存单是银行负债，承诺在约定到期日偿还本金，通常还支付利息。可转让存单可以在到期前转让，使初始投资者能够通过出售请求权获得流动性，而无需要求发行银行提前偿还。

从银行角度看，可转让存单是负债管理工具。它使银行能够按市场决定的利率筹资，而不必完全依赖零售客户存款。报价利率反映请求权期限、市场流动性、银行信用质量，以及其他短期投资的供给。

可交易性改变了投资者在到期前获取流动性的方式。不可转让的定期存款通常要求存款人持有至到期，或向银行申请提前支取；可转让存单则可在二级市场出售，但价格可能随市场利率及市场对发行银行信用质量的判断而变化。

\begin{table}[H]
\centering
\caption{定期存款与可转让存单}
\label{tab:ch7-cd-comparison}

\begin{tabularx}{\textwidth}{lXX}
\toprule
 & \textbf{不可转让定期存款} & \textbf{可转让存单} \\
\midrule

可转让性 & 通常不可转让 & 可在二级市场转让 \\

到期前流动性 & 通常需提前支取，或与银行另作安排 & 可通过出售请求权获得 \\

投资者风险暴露 & 主要面向发行银行 & 面向发行银行，并承担到期前的市场价格风险 \\

对银行的融资作用 & 存款融资 & 基于市场的批发融资 \\

\bottomrule
\end{tabularx}
\end{table}

中国银行间同业存单是市场化银行融资的重要例子。根据中国人民银行~2013~年暂行规定推出的此类记账式定期存款凭证，由符合条件的存款类机构在全国银行间市场发行。发行利率由市场决定，公开发行的存单可以在二级市场交易。它们为银行提供额外的批发融资来源，也为投资者提供对发行机构的可转让请求权。

\subsection{同业贷款与隔夜无担保融资}

银行日常代客户收付款，这些资金流动使各银行的流动性头寸暂时不同。同业拆借使资金盈余机构能够向存在短期缺口的机构提供资金。

无担保同业贷款是对借款银行的直接请求权。偿还不以特定抵押品作支持，贷款人依赖借款人的信用状况及双边或全市场信用限额。因此，利率同时反映一般流动性条件和交易对手信用风险。

隔夜期限在流动性快速再分配中具有核心作用。银行当天遇到非预期结算资金流出时，可以借入隔夜资金，并在次日资金流入后偿还。数天或数月期限的融资可在更长期间内锁定流动性，但也延长贷款人对借款人的风险暴露。

压力时期，请求权的无担保性质更加重要。银行若无法确定交易对手的信用状况，即使持有现金也可能拒绝放贷。许多机构采取类似行为时，市场可能在不同交易对手之间分割：流动性仍在体系内，却无法向需要它的机构重新分配。

信用限额与市场利率同样重要。如果贷款人不愿增加无担保风险暴露，即使报价利率大幅提高，也未必带来额外放贷。此时，可用信贷数量及交易对手身份与报价同样重要。

\subsection{回购协议与有抵押融资}

回购协议（repo）是一种短期融资契约，一方以证券换取现金，并约定未来按指定回购价格进行反向交易。从经济功能看，回购属于有抵押短期融资，但其法律结构因市场而异。

从现金借入方看，交易是正回购；从现金贷出方看，同一交易是逆回购。设市场价值为~$V_0$~的抵押品支持一笔回购，如果贷款人采用折扣率~$h$，则提供的现金为

\begin{equation}
C_0=(1-h)V_0.
\label{eq:ch7-repo-cash-haircut}
\end{equation}

折扣为抵押品价值下跌及借款人违约时的清算成本提供缓冲。如果回购期限为一年的~$\tau$，回购利率为~$r_{repo}$，则回购支付额为

\begin{equation}
C_T=C_0(1+r_{repo}\tau),
\label{eq:ch7-repo-repurchase-payment}
\end{equation}
这里对回购期间采用单利计息惯例。

\begin{example}{回购融资与折扣率}{repo-funding-haircut}

一家交易商持有市值人民币~1~亿元的政府证券，并通过一周期回购融资。现金贷款人采用~2\%~的折扣率。

提供的现金为

\[
C_0=(1-0.02)\times100=98
\]

百万元。如果年化回购利率为~2.5\%，市场采用每年~365~天的计息基础，则七天后的回购支付额为

\[
C_T=98\times\left(1+0.025\times\frac{7}{365}\right)
\approx 98.047
\]
百万元。

折扣率与回购利率发挥不同作用。回购利率是现金融资的价格；折扣率决定相对于抵押品价值能够获得多少现金。

\end{example}

抵押品能够减少贷款人的损失暴露。但抵押品价格可能下跌，证券可能难以清算，法律和操作问题也可能延迟贷款人取得抵押品。回购契约通常规定抵押品资格、估值程序、折扣率和保证金安排。

回购存续期内抵押品价值变化时，保证金调整可能要求补充现金或证券，以维持约定的保护水平。这种盯市机制保护贷款人，但抵押品价值下降时，可能给借款人带来流动性压力。

回购使交易商和其他机构能够融资债券及其他合格证券库存，是证券市场中介活动的核心安排。它也将货币市场与债券市场联系起来：证券价值与流动性影响其作为抵押品的适用性，回购融资的可得性与成本则影响中介持有该证券的意愿。

表~\ref{tab:ch7-money-market-instruments}~概括了本节所讨论工具的主要契约差异。

\begin{table}[H]
\centering
\caption{主要货币市场工具及其契约结构}
\label{tab:ch7-money-market-instruments}

\begin{tabularx}{\textwidth}{
@{}
>{\raggedright\arraybackslash
\hyphenpenalty=10000
\exhyphenpenalty=10000}p{2.3cm}
>{\raggedright\arraybackslash
\hyphenpenalty=10000
\exhyphenpenalty=10000}p{2.5cm}
>{\raggedright\arraybackslash
\hyphenpenalty=10000
\exhyphenpenalty=10000}p{2.6cm}
Y
@{}
}

\toprule

\textbf{工具}
&
\makecell[l]{\textbf{典型发行人}\\\textbf{或借款人}}
&
\makecell[l]{\textbf{抵押品}\\\textbf{或信用基础}}
&
\makecell[l]{\textbf{主要经济作用与风险}}
\\

\midrule

国库券 & 中央政府 & 主权支付义务；通常无担保 & 政府短期融资与流动性投资；优质主权市场中信用风险较低，到期前仍有市场价格风险 \\

商业票据 & 企业或金融机构 & 通常无担保；发行人信用 & 灵活的短期市场融资；发行人信用风险及滚动融资风险 \\

可转让存单 & 银行 & 银行存款支付义务；通常无担保 & 银行批发融资及可转让银行请求权；银行信用风险及滚动融资风险 \\

同业贷款 & 银行或其他金融机构 & 通常为无担保交易对手信用 & 短期流动性再分配；交易对手信用风险及滚动融资风险 \\

回购 & 交易商、银行、基金或其他参与者 & 以卖出并回购协议提供抵押 & 证券融资与流动性管理；抵押品估值、折扣率、交易对手及滚动融资风险 \\

资产支持商业票据 & 特殊目的载体 & 资产池，常附流动性或信用支持 & 资产支持的短期融资；基础资产质量、流动性及滚动融资风险 \\

\bottomrule
\end{tabularx}
\end{table}

\section{货币市场架构与定价}

货币市场价格形成于借款人、现金投资者、银行、交易商和投资基金之间的互动，并依托市场交易与结算基础设施。契约条款决定工具现金流，市场架构则影响融资获取、流动性再分配、价格发现与基准形成，以及交易执行与结算。

\subsection{银行、交易商与货币市场基金}

银行在货币市场中兼为借款人与贷款人。它们通过借款管理结算需求和资产负债表短期融资，并在出现临时流动性盈余时放贷。这些活动将短期融资市场与支付体系及更广泛的银行体系联系起来。

证券交易商将短期融资与债券等证券交易联系起来。持有证券库存的交易商可以通过回购融资部分头寸。这种资产负债表承载能力使交易商能够持有库存、为买卖双方提供即时交易服务，从而支持二级市场流动性。回购融资成本和可得性影响交易商的做市意愿。

货币市场基金汇集投资者现金，主要投资于短期限、高质量工具。它们为家庭、企业和机构提供分散化投资货币市场资产的渠道，同时将现金导向政府、银行、企业和交易商。在适用于基金及所在司法辖区的规则约束下，投资组合可包括政府短期证券、回购、商业票据和存单。

货币市场基金的结构使资产流动性与投资者流动性之间形成重要区别。投资者可能预期短时间内即可赎回份额，而部分组合资产在市场压力时期可能难以迅速出售。基金遭遇大量赎回时，可以动用现金和到期证券回款、出售资产，或减少新增购买。这些行为将投资者提款传导至短期借款人使用的市场。

\begin{figure}[H]
\centering
\begin{tikzpicture}[
  >=Latex,
  box/.style={
    draw,
    rounded corners=2pt,
    align=center,
    minimum height=1.0cm,
    text width=2.7cm,
    inner sep=5pt
  },
  widebox/.style={
    draw,
    rounded corners=2pt,
    align=center,
    minimum height=1.0cm,
    text width=3.6cm,
    inner sep=5pt
  },
  fundbox/.style={
    draw,
    rounded corners=2pt,
    align=center,
    minimum height=1.0cm,
    text width=1.5cm,
    inner sep=5pt
  },
  flow/.style={
    -Latex,
    line width=0.9pt,
    FinanceBlue
  },
  claim/.style={
    -Latex,
    dashed,
    line width=0.9pt,
    FinanceBlue
  },
abovelab/.style={
  font=\small,
  text=FinanceBlue,
  fill=white,
  align=center,
  fill=white,
  inner sep=1.5pt
},
  belowlab/.style={
    font=\small,
    text=black,
    fill=white,
    align=center,
    inner sep=1.5pt
  }
]

% =================================================
% Top row: Money market funds
% =================================================

\node[fundbox] (inv) at (-0.5,3.0)
  {现金投资者};

\node[fundbox] (mmf) at (3.8,3.0)
  {货币市场\\基金};

\node[widebox] (iss) at (9.2,3.0)
  {短期请求权发行人\\
   （政府、企业、银行）};

% Investor subscriptions and fund shares
\draw[flow]
  ([yshift=3pt]inv.east) --
  node[abovelab, midway, above=3pt] {现金\\认购}
  ([yshift=3pt]mmf.west);

\draw[claim]
  ([yshift=-3pt]mmf.west) --
  node[belowlab, midway, below=3pt] {基金份额}
  ([yshift=-3pt]inv.east);

% MMF investment in short-term claims
\draw[flow]
  ([yshift=3pt]mmf.east) --
  node[abovelab, midway, above=3pt] {现金\\投资}
  ([yshift=3pt]iss.west);

\draw[claim]
  ([yshift=-3pt]iss.west) --
  node[belowlab, midway, below=3pt]
  {国库券、\\商业票据、存单}
  ([yshift=-3pt]mmf.east);

% =================================================
% Middle row: Interbank lending
% =================================================

\node[box] (bank1) at (0,0)
  {流动性\\盈余银行};

\node[widebox] (bank2) at (9.2,0)
  {有临时融资\\需求的银行};

% Cash and loan claim
\draw[flow]
  ([yshift=3pt]bank1.east) --
  node[abovelab, midway, above=3pt]
  {同业资金}
  ([yshift=3pt]bank2.west);

\draw[claim]
  ([yshift=-3pt]bank2.west) --
  node[belowlab, midway, below=3pt]
  {贷款请求权}
  ([yshift=-3pt]bank1.east);

% =================================================
% Bottom row: Repo financing
% =================================================

\node[widebox] (cashl) at (0,-3.0)
  {现金贷出方\\（银行、货币基金、其他投资者）};

\node[widebox] (dealer) at (9.2,-3.0)
  {持有证券库存\\的交易商};

% 回购资金 and collateral
\draw[flow]
  ([yshift=3pt]cashl.east) --
  node[abovelab, midway, above=3pt]
  {回购资金}
  ([yshift=3pt]dealer.west);

\draw[claim]
  ([yshift=-3pt]dealer.west) --
  node[belowlab, midway, below=3pt]
  {证券抵押品}
  ([yshift=-3pt]cashl.east);

% =================================================
% Legend
% =================================================

\draw[flow] (0.5,-4.5) -- (1.4,-4.5);
\node[anchor=west, draw=none] at (1.5,-4.5)
  {\small 现金流};

\draw[claim] (4.2,-4.5) -- (5.1,-4.5);
\node[anchor=west, draw=none] at (5.2,-4.5)
  {\small 请求权或抵押品转移};

\end{tikzpicture}

\caption{货币市场融资的机构渠道}
\label{fig:ch7-money-market-channels}
\end{figure}

该图为示意图。银行可同时处于市场借贷两侧，货币市场基金与交易商则履行不同的中介功能。图示旨在说明短期资金如何通过多种机构渠道流动。

\subsection{基准利率、利差与价格发现}

中央银行通过结算余额供给、公开市场操作和流动性便利影响货币市场条件。这些操作影响极短期利率及总体融资可得性。但交易对手信用风险、抵押品条件、市场分割和机构特有融资条件，仍可能使各工具利率不同。

货币市场契约常参照基准利率定价。有效基准应代表界定清楚的市场区间，依据可靠信息，并有透明方法和适当治理支持。其经济含义取决于基础交易及利率构造方法。

基准在所衡量的融资类型和构造方式上均存在差异。一项区别是融资是否有担保。有担保利率反映抵押品支持的融资交易，受抵押品质量和可得性影响；无担保利率反映无抵押品借款，更直接体现交易对手信用条件。另一项区别是利率形成方式：交易型基准由合格市场交易推导，报价型基准依赖选定市场参与者提交的利率。

中国货币市场体现了这些差异。上海银行间同业拆放利率（Shibor）是报价型无担保基准，由报价银行团提交利率形成，并发布多个期限品种。相比之下，存款类机构间回购利率（DR）是交易型有担保利率，依据存款类机构之间以利率债为抵押品的质押式回购交易形成。七天期利率~DR007~广泛用作银行体系流动性指标和货币市场定价参考。

国际基准改革增加了交易型隔夜利率的使用。美国~SOFR~基于有担保的国债回购交易，英国~SONIA~和欧元区~€STR~则基于无担保隔夜交易。

表~\ref{tab:ch7-benchmark-comparison}~概括了这些基准的主要特征。

\begin{table}[H]
\centering
\caption{主要货币市场基准及其特征}
\label{tab:ch7-benchmark-comparison}

\begin{tabularx}{\textwidth}{
@{}
>{\raggedright\arraybackslash}p{1.7cm}
>{\raggedright\arraybackslash}p{3.2cm}
>{\raggedright\arraybackslash}p{2.0cm}
Y
@{}
}

\toprule

\textbf{基准}
&
\textbf{基础市场}
&
\textbf{融资基础}
&
\textbf{利率形成方式}
\\

\midrule

Shibor & 人民币同业融资 & 无担保 & 银行团报价 \\

DR007 & 七天期同业质押式回购 & 有担保 & 实际交易加权利率 \\

SOFR & 隔夜美国国债回购 & 有担保 & 基于交易的成交量加权中位数 \\

SONIA & 隔夜英镑存款 & 无担保 & 基于交易的截尾均值 \\

€STR & 隔夜欧元批发融资 & 无担保 & 基于交易计算 \\

\bottomrule
\end{tabularx}
\end{table}

对于短期工具及期限和市场特征可比的基准，利差为

\begin{equation}
s_i=r_i-r_b.
\label{eq:ch7-money-market-spread}
\end{equation}
其中，$r_i$~表示短期工具利率，$r_b$~表示适当的基准利率。

利差扩大可能反映信用风险增加、市场流动性下降、资产负债表约束收紧，或工具需求减弱。因此，同一观察到的利差可能来自不同基础条件。

借款人与贷款人根据融资需求、新信息、货币操作和交易条件调整利率，形成价格发现。货币市场请求权期限短、重新定价频繁，使极短期利率对融资条件变化尤其敏感。纳税或报表日期前后、央行流动性变化后，或机构放贷意愿下降时，隔夜利率可能剧烈波动。

基准利率也用于货币市场之外的契约。浮动利率贷款、债券、互换及衍生品可能参照短期基准确定契约支付。因此，货币市场价格发现影响更广泛金融请求权的定价与估值。

\subsection{抵押品、清算与结算}

除借款人的一般信用状况外，抵押品为贷款人提供额外回收来源。其有效性取决于三个问题：哪些资产合格、如何估值，以及违约后贷款人能多快行使相关权利。

政府证券及其他高质量有价证券常用作货币市场抵押品。其价值较为透明，通常比流动性较低的资产更易出售或用于融资。其他证券也可能合格，但在价格波动、信用风险或清算成本较高时，贷款人可能采用更高折扣率。

初始折扣为抵押品价值下跌提供缓冲，但市场价格变化后，缓冲可能不足。因此，回购契约可能要求每日或日内对抵押品盯市。如果抵押品价值低于协议要求，借款人可能需要补充证券或现金。这类追缴保证金要求限制贷款人风险暴露，但也可能给借款人带来流动性压力，抵押品价格急跌时尤其如此。

清算确定交易产生的义务，并在适用情况下确定待结算金额。结算通过转移现金和证券完成义务。证券交易中，卖方须交付证券，买方须支付相应现金。如果两项转移发生在不同时点，任何一方都可能面临对方不履约的风险。券款对付（DvP）将证券交付与相应现金支付相联系，降低这种本金风险。

中国银行间市场体现了交易服务与证券登记、托管和结算之间的分工。中国外汇交易中心（CFETS），亦称全国银行间同业拆借中心，为同业存单提供发行、交易和信息服务；上海清算所（SHCH）负责登记、托管和结算。这一安排将工具交易与负责维护所有权记录、完成结算的机构职能区分开来。

这些安排无法消除流动性风险。结算延误或操作问题妨碍抵押品动用时，即使有价值的抵押品也可能无法及时提供融资。货币市场流动性因此不仅取决于金融请求权质量，也取决于其转移所依托的基础设施。

\begin{tcolorbox}[
  enhanced,
  breakable,
  colback=white,
  colframe=FinanceBlue,
  boxrule=0.7pt,
  sharp corners,
  borderline west={2.2pt}{0pt}{FinanceBlue},
  title=\textbf{金融视角：抵押品必须能够动用},
  coltitle=white,
  colbacktitle=FinanceBlue,
  fonttitle=\bfseries,
  before skip=1.2em,
  after skip=1.2em
]

抵押品价值不仅取决于信用质量和市场价格，还取决于融资需要出现时，能否转移、被接受和清算。法律、操作与结算摩擦可能限制其作为即时流动性来源的作用。

\end{tcolorbox}

\section{货币市场的流动性压力}

正常条件下，货币市场可能表现出较高流动性。短期限和定期再融资使资金能够在机构间顺畅流动。但贷款人不愿续借时，同样的特征也可能加剧压力。大量债务可能短时间内到期，迫使借款人在长期资产能够出售或到期之前筹集现金。

\subsection{融资挤兑与市场冻结}

融资挤兑指短期债权人撤回融资、拒绝续借到期请求权，或缩短新贷款期限。即使机构依赖短期批发融资而非存款，其动态过程也类似银行挤兑。市场冻结则指参与者不愿或无法按现行条款交易，导致交易或新增融资急剧下降。

假设借款人以每日或每周到期的债务融资长期资产。只要贷款人继续续借，就能维持头寸。对借款人资产质量、信用状况的担忧，或贷款人自身流动性需求，可能降低续借意愿。借款人必须在现有请求权到期时筹集现金。

为履行义务，借款人可动用承诺授信、使用现金储备、追加抵押品，或出售资产。如果许多机构面临相似融资压力，这些反应可能相互强化。资产出售可能压低市场价格，流动性需求上升则进一步挤压融资条件。资产价格下跌还可能削弱资产负债表、降低抵押品价值，使融资压力进一步增大。这一过程可能自我强化：

\begin{figure}[H]
\centering
\begin{tikzpicture}[
  >=Latex,
  box/.style={draw, rounded corners, align=center, text width=4.0cm, minimum height=0.9cm},
  flow/.style={-Latex, line width=0.8pt, FinanceBlue}
]

\node[box] (concern) at (0,3.2) {对信用、流动性\\或抵押品的担忧};
\node[box] (roll) at (6.0,3.2) {贷款人缩短期限\\或拒绝续借};
\node[box] (cash) at (6.0,0.7) {借款人筹集现金\\并出售资产};
\node[box] (price) at (0,0.7) 
{资产价格下降，削弱\\抵押品价值与市场流动性};

\draw[flow] (concern) -- (roll);
\draw[flow] (roll) -- (cash);
\draw[flow] (cash) -- (price);
\draw[flow] (price) -- (concern);

\end{tikzpicture}
\caption{短期融资的反馈循环}
\label{fig:ch7-funding-run-loop}
\end{figure}

在无担保市场中，即使总体流动性仍充裕，对借款人质量的不确定性也可能减少放贷。机构可能更愿保留流动性，而非增加对状况难以评估的交易对手的风险暴露。融资压力同时影响价格与数量：利差可能上升，部分借款人则可能完全丧失融资渠道。

商业票据市场与货币市场基金也可能出现类似压力。投资者可能从商业票据转向政府证券或现金，减少发行人可获得的融资，增加其对承诺银行信贷额度的依赖。如果货币市场基金减少购买商业票据与存单，或出售资产应对提款，快速赎回就会强化这一过程。压力由此从现金投资者传向短期借款人，再传向银行资产负债表。

2013~年~6~月中国银行间市场的情况说明了这些机制。随着融资不确定性增强，机构更不愿放贷，削弱了银行体系内部的流动性再分配。短期利率急升，随后在流动性条件与预期稳定后回落。

\subsection{回购折扣率、抵押品压力与避险转移}

回购融资条件同时取决于回购利率与折扣率。回购利率决定现金融资成本，折扣率决定给定数量抵押品能够借入多少现金。

市场压力时期，这一区别尤其重要。考虑完全由回购借款与交易商自有资金融资的简化交易商头寸。设~$V$~为证券市值，$E$~为交易商投入的自有资金。按标准折扣惯例，回购折扣率~$h$~表示贷款人针对价值~$V$~的抵押品提供~$(1-h)V$~现金。因此，交易商必须以自有资金承担剩余部分：

\begin{equation*}
E=hV.
\label{eq:ch7-haircut-equity}
\end{equation*}

在这一简化融资结构下，资产权益比为

\begin{equation*}
\frac{V}{E}=\frac{1}{h}.
\label{eq:ch7-haircut-leverage}
\end{equation*}

市场压力时期折扣率上升，其融资影响尤为显著。

\begin{example}{折扣率上升产生融资需求}{haircut-funding-need}

一家交易商通过回购融资市值人民币~1~亿元的证券。折扣率为~2\%~时，贷款人提供~9800~万元，交易商投入~200~万元自有资金。

如果抵押品价值仍为~1~亿元，而折扣率升至~5\%，回购融资降至~9500~万元。交易商必须额外提供~300~万元现金或抵押品，才能维持相同头寸。

如果许多交易商面临相同变化，它们可能缩减头寸或出售证券释放现金。保护单个贷款人的折扣率，在全市场同时上升时，也可能推动整个体系去杠杆。

\end{example}

折扣率提高与抵押品价值下降同时发生时，抵押品压力加剧。借款人既面临每单位抵押品借款能力下降，也面临现有抵押品价值缩水。追缴保证金可能迫使其在下跌市场中出售资产，强化融资流动性与市场流动性的相互作用。

抵押品质量还影响流动性流向。压力时期，贷款人往往更愿接受政府证券，而较少愿意接受私人部门或流动性较低的抵押品。投资者同样将现金转向被视为更安全、更有流动性的请求权。这种\emph{避险转移}可能降低高质量短期政府请求权的收益率，同时扩大私人部门货币市场工具的利差。它体现流动性在资产与交易对手之间的重新配置，未必意味着总体流动性下降。

抵押品、折扣率、交易商资产负债表与市场流动性的相互作用，解释了回购市场为何影响其直接参与者之外的市场。如果交易商融资能力下降，其为债券市场提供中介服务的能力与意愿可能减弱。市场流动性下降又可能增加价格波动，降低用作抵押品的证券价值。贷款人可能进一步提高折扣率，从而收紧融资条件。
\begin{tcolorbox}[
  enhanced,
  breakable,
  colback=white,
  colframe=FinanceBlue,
  boxrule=0.7pt,
  sharp corners,
  borderline west={2.2pt}{0pt}{FinanceBlue},
  title=\textbf{金融视角：利率不变时，更高折扣率也能收紧融资},
  coltitle=white,
  colbacktitle=FinanceBlue,
  fonttitle=\bfseries,
  before skip=1.2em,
  after skip=1.2em
]

回购利率不变时，融资条件仍可能显著收紧。更高折扣率减少给定数量抵押品能够筹集的现金。借款人须投入更多自有资金、追加抵押品，或缩减头寸。因此，即使回购报价利率变化不大，折扣率也会影响可获得的融资数量。

\end{tcolorbox}

%========================================================
% Summary
%========================================================

\TailModule{summary}{本章小结}

货币市场重新分配流动性，并为短期现金流错配提供融资。如果现金流入与流出时点不同，即使基础资产和项目具有正的长期价值，借款人也可能需要短期资金。短期融资减少持有大量闲置现金的需要，但反复再融资产生滚动融资风险。

货币市场工具的契约结构不同。国库券是短期主权请求权；商业票据通常是无担保企业或金融债务；可转让存单是可转让银行负债；同业贷款在金融机构间提供无担保融资；回购则以证券为支持提供现金。报价惯例十分重要，贴现工具尤其如此；还应区分回购利率与抵押品折扣率。

货币市场架构影响定价与流动性。银行兼任借贷双方，交易商融资证券库存，货币市场基金则将投资者现金导向短期请求权。基准利率因融资是否有担保、依据交易还是报价而异。抵押品资格与估值，以及清算结算安排，共同影响使用短期请求权和证券获取流动性的能力。

贷款人拒绝续借短期融资，或大量赎回减少对货币市场借款人的资金供给时，可能出现流动性压力。回购市场中，更高折扣率减少抵押品支持的可融资金额。如果折扣率上升与抵押品价格下跌同时发生，追缴保证金和资产出售可能强化融资流动性与市场流动性的相互作用。短期限在正常条件下提供灵活性，但再融资不确定时，也可能增加脆弱性。

%========================================================
% Key Terms
%========================================================

\TailModule{keyterms}{关键术语}

货币市场（money market）；短期融资（short-term finance）；流动性（liquidity）；现金流时点（cash-flow timing）；营运资本（working capital）；财务灵活性（financial flexibility）；滚动融资风险（rollover risk）；国库券（Treasury bill）；持有期回报率（holding-period return）；贴现收益率（discount yield）；商业票据（commercial paper）；可转让存单（negotiable certificate of deposit）；同业贷款（interbank loan）；回购协议（repurchase agreement）；逆回购（reverse repo）；回购利率（repo rate）；抵押品（collateral）；折扣率（haircut）；追缴保证金（margin call）；货币市场基金（money market fund）；基准利率（benchmark rate）；交易型基准（transaction-based benchmark）；报价型基准（quotation-based benchmark）；上海银行间同业拆放利率（Shibor）；DR007；价格发现（price discovery）；券款对付（delivery versus payment）；融资挤兑（funding run）；市场冻结（market freeze）；避险转移（flight to quality）；融资流动性（funding liquidity）；市场流动性（market liquidity）。

%========================================================
% Financial AI Lab
%========================================================

\TailModule{aiprompt}{金融~AI~实验室}

\textbf{实验：重建回购交易并进行压力测试}

请~AI~系统构造一笔一周期回购，使证券交易商融资债券头寸。要求回答列明抵押品市值、折扣率、提供的现金、回购利率、日计数惯例及回购支付额。

独立审核交易。核实折扣率用于抵押品价值并决定提供的现金金额，回购利率则决定融资成本。同时检查~AI~是否正确区分正回购与逆回购视角。

随后加入两项冲击：抵押品价值下降~4\%，折扣率提高~3~个百分点。保持原有回购融资金额不变，请~AI~计算产生的融资缺口，以及恢复原融资头寸所需的额外现金或合格抵押品。用本章契约关系核验结果。

最后，请~AI~解释折扣率上升如何推动全市场去杠杆。检查回答是否错误地将抵押品视为消除了流动性风险，或假设抵押品总能立即按其市场报价出售或质押。

%========================================================
% Discussion Questions
%========================================================

\TailModule{questions}{讨论题}

\begin{enumerate}

\item 为什么盈利且有偿付能力的企业仍可能需要短期融资？区分流动性问题与偿付能力问题。

\item 比较商业票据与回购协议。有无抵押品如何影响贷款人风险暴露、融资条款和借款人的滚动融资风险？

\item 为什么即使银行体系总体流动性看似充足，无担保同业市场仍可能缺乏流动性？

\item 为什么无担保货币市场利率与有担保回购利率之间的利差可能因不同原因扩大？解释这种变化需要哪些信息？

\end{enumerate}

%========================================================
% Applied Problems
%========================================================

\TailModule{appprobs}{应用题}

\begin{enumerate}

\item \textbf{现金流时点与短期融资。}

一家企业月初持有现金人民币~400~万元。该月预计收到现金~1800~万元、支付~2700~万元。

\begin{enumerate}

\item 计算新增融资前的净现金头寸。

\item 计算短期融资需求。

\item 假设原定本月收回的~700~万元应收账款推迟至下月收回，重新计算融资需求。

\item 解释为什么这一变化本身不足以证明企业资不抵债。

\end{enumerate}

\item \textbf{货币市场报价惯例。}

一张~90~天期票据，面值为~1000，售价为~988。采用每年~360~天的计息基础。

\begin{enumerate}

\item 计算持有期回报率。

\item 计算以购买价格为基础的年化货币市场收益率。

\item 计算以面值为基础的贴现收益率。

\item 解释为什么基础现金流相同，两种年化收益率仍不同。

\end{enumerate}

\item \textbf{商业票据利差。}

一家高评级企业以~3.8\%~的年化利率发行~60~天期商业票据，此时可比短期基准利率为~2.9\%。

\begin{enumerate}

\item 计算相对于基准的利差。

\item 假设基准升至~3.4\%，利差扩大至~1.6~个百分点，新的商业票据利率为多少？

\item 将融资利率增幅分解为基准成分和利差成分。

\item 给出利差可能扩大的两个经济原因。

\end{enumerate}

\item \textbf{回购利率与折扣率。}

一家交易商持有价值人民币~2.5~亿元的证券，并以隔夜回购融资。初始折扣率为~3\%，年化回购利率为~2.2\%，按每年~365~天计息。

\begin{enumerate}

\item 计算提供的现金金额。

\item 计算次日的回购支付额。

\item 如果抵押品价值不变、折扣率升至~6\%，交易商需额外投入多少自有资金，才能维持相同资产头寸？

\item 解释为什么折扣率提高与回购利率提高相同百分点，具有不同经济影响。

\end{enumerate}

\end{enumerate}

%========================================================
% Data Application
%========================================================

\TailModule{dataapp}{数据应用}

\textbf{中国货币市场的流动性条件}

从中国货币网或中国人民银行等权威来源，获取至少两种中国短期货币市场利率的日度数据。一组可用的组合为~DR007~与七天期银行间回购利率或七天期~Shibor。如果加入央行七天期公开市场操作利率，应将其与市场决定的利率明确区分。

选择同时包括正常时期和至少一次明显短期波动的样本期间，绘制所选利率。计算市场利率之间的利差，识别利差或利率水平急剧变化的日期。

结合基础市场结构解释变化。明确各利率是否有担保、哪些机构参与、哪些抵押品合格，以及基准依据交易还是报价。检查月末、季末、缴税、节假日或其他流动性管理日期，是否与异常波动重合。

不要仅将利差变化归因于信用风险或货币政策。解释需要哪些额外证据，才能区分总体流动性、交易对手风险、抵押品条件和市场分割的变化。报告样本期间、利率定义、计算惯例及原始数据来源。

%========================================================
% Case Analysis
%========================================================

\TailModule{caseanalysis}{案例分析}

\textbf{案例：2020~年~3~月美国短期融资市场的压力}

2020~年~3~月，新冠疫情迅速蔓延，引发对现金和高流动性资产的异常需求。优质型货币市场基金（prime money market funds），尤其是面向机构投资者的此类基金，遭遇大量赎回；商业票据及其他私人部门短期工具市场承受严重压力。基金应对赎回时，减少持有和购买私人部门短期请求权，使企业和金融机构为到期义务再融资更加困难。

美联储通过多项紧急便利作出应对。2020~年~3~月~17~日宣布的商业票据融资便利（CPFF），通过特殊目的载体购买合格商业票据，提供流动性后盾。3~月~18~日宣布的货币市场共同基金流动性便利（MMLF），向合格金融机构提供贷款，以其从货币市场基金购买的合格资产作担保。这些及其他干预宣布并实施后，短期融资市场压力缓解。

运用本章框架分析该事件，追踪从货币市场基金投资者到商业票据发行人的融资链条。解释即使发行人经营现金流尚未立即恶化，投资者赎回为什么仍可能减少其可获得的融资。区分市场流动性与融资流动性，并识别滚动融资风险在哪些环节进入链条。

随后比较~CPFF~与~MMLF~的作用。各项便利针对什么市场摩擦？即使后续实际使用有限，可信后盾的存在为何仍可能改善市场运行？如果参与者预期未来短期融资压力时期会获得公共流动性支持，可能产生哪些激励问题？

\emph{原始资料来源：美国联邦储备理事会，“商业票据融资便利”，2020~年~3~月；美国联邦储备理事会，“货币市场共同基金流动性便利”，2020~年~3~月。}

%========================================================
% Further Reading
%========================================================

\TailModule{furtherread}{延伸阅读}

\begin{itemize}

\item Stigum, Marcia L., and Anthony Crescenzi (2007),
\emph{Stigum's Money Market（斯蒂格姆货币市场）}, 4th ed., McGraw-Hill.

\item Gorton, Gary, and Andrew Metrick (2012),
``Securitized Banking and the Run on Repo（证券化银行与回购挤兑）,''
\emph{Journal of Financial Economics}, 104(3), 425--451.

\item Covitz, Daniel, Nellie Liang, and Gustavo A. Suarez (2013),
``The Evolution of a Financial Crisis: Collapse of the Asset-Backed Commercial Paper Market（金融危机的演进：资产支持商业票据市场的崩溃）,''
\emph{The Journal of Finance}, 68(3), 815--848.

\item Krishnamurthy, Arvind, Stefan Nagel, and Dmitry Orlov (2014),
``Sizing Up Repo（衡量回购市场规模）,''
\emph{The Journal of Finance}, 69(6), 2381--2417.

\item Committee on the Global Financial System (2017),
\emph{Repo Market Functioning（回购市场运行）},
CGFS Papers No.~59, Bank for International Settlements.

\end{itemize}

\ChapterTail{summary,keyterms,aiprompt,questions,appprobs,dataapp,caseanalysis,furtherread}
